% !TEX TS-program = lualatexmk
% !TEX parameter =  --shell-escape
\documentclass[12pt]{book}

%%%%%%%%%%%%%%% STANDARD PACKAGES %%%%%%%%%%%%%%%

%% for personal reasons, the commands are a mix of French and English names. Feel free to choose one or the other if you make use of them. This file is subject to change later on.

\usepackage{geometry}
\usepackage{fontspec}
\usepackage[english,ecclesiasticlatin.usej,activeacute]{babel}
     \babelprovide[hyphenrules=latin]{ecclesiasticlatin}
%          \usepackage{xspace} %% better to use {} after certain macros (or even within definitions) 
     \usepackage{multicol}
\usepackage{paracol}
\footnotelayout{m} %% this allows merged footnote if text is in parallel columns (mostly translations)
\usepackage{fancyhdr}
\usepackage{needspace} %%pour réserver de l'espace en bas de page ; ça évite des séparations entre les titres et la partition ou le texte qui les suivent.
\usepackage[compact]{titlesec}
\usepackage{emptypage}
\usepackage{xcolor}
\usepackage{perpage}%%%to reset footnotes at each page%%%
%\usepackage{xstring}
\usepackage{enumitem} %% nécessaire pour les textes des psaumes.
\usepackage{etoolbox}
\usepackage{lettrine}
%\usepackage{tocloft}%%to fix toc title. titlesec companion packages could also work.
\usepackage{indentfirst}%% to allow for \section with indented 1st paragraphs in good Franco-Latin style (LaTeX default is to suppress this. Package sets  \@afterindentfalse to (always) true.
\usepackage[savepos]{zref}
\usepackage{hyperref}
\usepackage{refcount}


%%%%%%%%%%%%%%% HYPHENATION AND TYPOGRAPHICAL CONVENTIONS %%%%%%%%%%%%%%

%% page settings %%%%
\sloppy %% to avoid overfull hboxes

%\setlength{\textwidth}{5.75in}

\geometry{letterpaper}
\geometry{bindingoffset=0.5in,
inner=1in,
outer=1.35in,
top=1.25in,
bottom=1.25in,
headsep=0.75in
} %%

%% General scale of all graphical elements. Also borrowed from MB
%% Values different from 1 are largely untested.
%% Used in those commands (e.g. everything FontSpec) that use a scale parameter.
\newcommand{\customscale}{1}

\AtBeginDocument{\setlength{\parindent}{1em}} %% should set 1em to EB Garamond not Computer Modern.

%%%%%%%%%%%%%%% GREGORIO CONFIG %%%%%%%%%%%%%%%

\usepackage[autocompile]{gregoriotex}

%% sets * and † to font called via fontspec
\def\GreStar{*}
\def\GreDagger{†}

%%prints asterisk instead of star normally inserted by <v>\greheightstar</v>
 \gredefsymbol{greheightstar}{EB Garamond}{asterisk}

%% should prevent flat turning custos into a clef, as Solesmes  uses an altered custos in only one chant per Élie Roux.
\gresetcustosalteration{invisible}

%produces a score with a smaller initial (default is 40 pt) and annotations like in the Solesmes books; the 1st is optional and can be omitted as needed.
 \NewDocumentCommand{\gscore}{O{}mm}{%
 \greannotation{#1}%
\greannotation{#2}%
\grechangestyle{initial}{\fontsize{28}{28}\selectfont}%
 \gregorioscore{partitions/#3}}
 
 
  \NewDocumentCommand{\lectionscore}{O{}m}{%
   \grecommentary{#1}%
  \grechangestyle{initial}{\fontsize{36}{36}\selectfont}%
  \gregorioscore{partitions/#2}}
 
 %% default should be 40, if not using, annotations need to be turned off in gabc. For the first antiphon of the office, Gospel canticle etc. (initial is both to refer to first and the large illustrated initial, used sometimes with an annotation, cf. Christus factus est of Tenebrae in the ritus monasticus.
  \NewDocumentCommand{\initialscore}{O{}m}{%
      \grechangedim{initialraise}{0.5cm}{scalable}%
  \greannotation{#1}%
   \grechangestyle{initial}{\fontspec{EB Garamond Initials}\fontsize{45}{45}\selectfont}%
 \gregorioscore{partitions/#2}
   \grechangedim{initialraise}{0cm}{scalable}}

%% score with no initial, e.g psalms, common tones, etc.
\newcommand{\smallscore}[2][y]{
  \gresetinitiallines{0}
  \gregorioscore{partitions/#2} %% à remplacer avec NewDocumentCommand ; les arguments ne marchent pas correctement … que fait le [y] ?
  \gresetinitiallines{1}
}

\NewDocumentCommand{\zerobaroffsettextleft}{}%
{%
\grechangedim{maxbaroffsettextleft}%
{0cm}%
{scalable}%
}

\NewDocumentCommand{\zerobaroffsettextright}{}%
{%
\grechangedim{maxbaroffsettextright}%
{0cm}%
{scalable}%
}

%%%the default for left bar offset is 0.3cm. Right is 0.15cm. This needs to be applied after a score where the bar offset is modified.

\NewDocumentCommand{\resetbaroffsettextleft}{}%
{%
\grechangedim{maxbaroffsettextleft}%
{0.3cm}%
{scalable}%
}

\NewDocumentCommand{\resetbaroffsettextright}{}%
{%
\grechangedim{maxbaroffsettextright}%
{0.15cm}%
{scalable}%
}

 \NewDocumentCommand{\MagAnnotation}{}{%
 \grechangedim{annotationseparation}{-0.01cm}{fixed}%
 } %% pour les antiennes où le chiffre convient mieux aux miniscules (parfois 1, 2, 4, 7…) ; il n'est pas nécessaire pour 6 ou 8.
%\gresetheadercapture{annotation}{greannotation}{string}

%% A-bove L-ines T-ext shall be italicized and in a smaller font
\grechangestyle{abovelinestext}{\itshape\fontsize{8}{10}\selectfont}
%% for psalm tones etc. Roman/upright text is what the Liber Usualis uses. 
\NewDocumentCommand{\altnormal}{}{\grechangestyle{abovelinestext}{\normalfont\fontsize{8}{10}\selectfont}}
%%  resets ALT font
\NewDocumentCommand{\altitshape}{}{\grechangestyle{abovelinestext}{\itshape\fontsize{8}{10}\selectfont}} 

\NewDocumentCommand{\altlarge}{}{\grechangestyle{abovelinestext}{\large}} %% for I reading of Lamentations

\NewDocumentCommand{\altupright}{}{\grechangestyle{abovelinestext}{\normalfont}} %%


%% fine-tuning of space beween the staff and the text above lines (used for Magnificats; needs to be rethought (again, should 2 and 3 be raised or not?) and worked into \smallscore
\newcommand{\altraise}{0.1cm} %% default is -0.1cm %% needs to be fine-tuned for \rubrique command with EB Garamond font. -0.2cm is MB value
\grechangedim{abovelinestextraise}{\altraise}{scalable}

\newcommand{\altheight}{0.5cm} %% default is 0.3cm and value must be bigger than text height; this should fix the problem. but needs further investigation.
\grechangedim{abovelinestextheight}{\altheight}{scalable} %% this is a very finicky command.

\newcommand{\comraise}{0.4cm} %% default is 0,2m %%
\grechangedim{commentaryraise}{\comraise}{scalable}

%% allows printing of glyphs from Gregorio score font (available glyphs listed in documentation) as text.
\makeatletter
\def\gretextglyph#1{{\gre@font@music\csname GreCP#1\endcsname}}
\makeatother

%%% Font %%%
\setmainfont{EB Garamond}[UprightFont=EB Garamond Regular,
ItalicFont= EB Garamond Italic,
BoldFont= EB Garamond Bold,
Ligatures=Rare,
Numbers=Proportional,
Numbers=OldStyle] %% all \kern numbers are based on this and can be removed if this font is abandoned.
\newfontfamily\symbolfont{liturgy}
  %% text must be written {\symbolfont{+}} (you can use V and R as well) but may conflict as + is † in gabc. Cross should be \small or even \footnotesize


%    \grechangestyle{initial}{\fontspec{EB Garamond Initials}\fontsize{#1}{#1}\selectfont}

%%% default initial size is 32 points
%\newcommand{\defaultinitialsize}{28}
%\initialsize{\defaultinitialsize}

%%%% Font %%%
%\setmainfont{EB Garamond}[UprightFont=EB Garamond Regular,
%ItalicFont= EB Garamond Italic,
%BoldFont= EB Garamond Bold,
%Ligatures=Rare,
%StylisticSet=6,
%Numbers=Proportional,
%Numbers=OldStyle] %% all \kern numbers are based on this and can be removed if this font is abandoned. %%StylisticSet=6 is, in Pardo EBGaramond, the long-tailed Q.
%\newfontfamily\symbolfont{liturgy} 
  %% text must be written {\symbolfont{+}} (you can use V and R as well) but may conflict as + is † in gabc. Cross should be \small or even \footnotesize
  
  %%%Roman numerals for the year %% https://tex.stackexchange.com/questions/185548/the-year-in-roman-and-the-month-in-text

\makeatletter
\newcommand{\YEAR}{\@Roman{\the\year}}
\makeatother


%%% TYPOGRAPHICAL AND SYMBOL COMMANDS  %%%%

%% \P is pilcrow and is defined in LaTeX as is.
%% Add {\textcolor{gregoriocolor} as first part of command's definition if changing to red.


%%for proper spacing of all-caps letters to prevent clashes, e.g. serif strokes touching.

\NewDocumentCommand\capspace{m}{{\addfontfeature{LetterSpace=5.0}{#1}}}
\NewDocumentCommand\scspace{m}{\textsc{{{\addfontfeature{LetterSpace=5.0}{#1}}}}}

\NewDocumentCommand\feasttitle{m}{{\centering{\capspace{\huge{#1}}}\par}}


%% macro to print Alleluia for versicles.
\NewDocumentCommand{\tpalleluia}{}{(\textit{T.P.} Allelúia.)}

\newcommand{\specialcharhsep}{3mm} % space after invoking R/ or V/ or A/ outside rubrics
%\newcommand{\vv}{\textcolor{gregoriocolor}{\fontspec[Scale=\customscale]{Charis SIL}℣.\nolinebreak[4]\hspace{\specialcharhsep}\nolinebreak[4]}} %% format from MB

\newcommand{\vv}{{\normalfont ℣.\nolinebreak[4]\hspace{\specialcharhsep}\nolinebreak[4]}}
\newcommand{\rr}{{\normalfont ℟.\nolinebreak[4]\hspace{\specialcharhsep}\nolinebreak[4]}}

%% typesets a cross pattée.
\NewDocumentCommand{\cc}{}{%
    % Grouping to keep font changes local
    {%
        % Ensure we're not in italics (since liturgy.ttf doesn't have italic)
        \normalfont%
        % Select the font liturgy.ttf
        \symbolfont%
        % Set the size
        \footnotesize%
        % In liturgy.ttf, the plus (+) is a cross pattée
        +%
    }%    
} %%  can replace <+> with <U+2720> (see above) if a suitable font is found (or EB Garamond font is fixed); command will be useful in lieu of typing the Unicode.
%% use <v> tag to insert into a gabc score (usually for the faithful or for blessings, e.g. the font.

%% Same special characters, for in-score use (<sp>V/ R/ A/ +</sp>)

\gresetspecial{V/}{{\fontspec[Scale=\customscale]{EB Garamond}℣.~}}
\gresetspecial{R/}{{\fontspec[Scale=\customscale]{EB Garamond}℟.~}}
\gresetspecial{+}{{\fontspec[Scale=\customscale]{EB Garamond}†~}}
\gresetspecial{*}{\gresixstar}
\gresetspecial{cross}{\cc}

%% Same special characters, for use in rubrics (no space)

\newcommand{\vvrub}{{\normalfont ℣.~}}
\newcommand{\rrrub}{{\normalfont ℟.~}}

\NewDocumentCommand{\antiphona}{mm}
{\begin{otherlanguage*}{#1}%
\noindent%
 Ant. #2.%
 \end{otherlanguage*}%
 }
 


%%rubrics: black italics, smaller than body of psalms etc

\NewDocumentCommand{\rubrique}{m}{%
    {%
        \fontsize{10}{12}%
        \selectfont%
        \textit{%
        %
        #1%
    }%
    }}
    
%% macro to print normal text inside of rubric (name of a chant or prayer, etc.)

\NewDocumentCommand{\normaltext}{m}{%
    {%
        \normalfont%
        #1%
    }%
    }
    
%% in case something should be bolded inside of a rubrique
\NewDocumentCommand{\rubriquegras}{m}{%
    {%
        \normalfont%
\textbf{#1%
}%
}}

%% to print in red instead of italicizing
\NewDocumentCommand{\rouge}{m}{%
\textcolor{gregoriocolor}%
{\fontsize{10}{12}%
\selectfont{%
#1%
}%
}}

%% to print in black within rouge group.

%\NewDocumentCommand{\textenoir}{m}{%

\NewDocumentCommand{\textenoir}{m}{%
\textcolor{black}%
{%
\normalfont%
#1%
}%
}

%% rouge but in italic (this is technically not correct).
\NewDocumentCommand{\rougeit}{m}{%
\textcolor{gregoriocolor}%
{\fontsize{10}{12}%
\selectfont%
\textit{#1%
}%
}}

%%for Scriptural references of the chapter.
\NewDocumentCommand{\scripture}{m}{%
{\raggedleft{\textit{#1}%
}%
\par}%
}

%\newcommand{\rubriquegras}[1]{{\fontsize{9}{11}\selectfont\textbf{#1}}}

%%macro to space punctuation with ecclesiasticlatin language from babel.
\NewDocumentCommand{\espaceponctuation}{}{\hspace{0.10em}} %%this value seems more balanced with this font than the definition provided in the ecclesiasticlatin documentation.

\NewDocumentCommand{\myrule}{}{%
    \par%
    {%
        \centering%
        \rule{0.3\textwidth}{0.4pt}%
        \par%
    }% typesets a horizontal rule like on the page with the prayers before and after the office.
}  %%If you're using color at all in your document, you might want to either force \myrule to use black or make it cusotmizable. (from u/Independent-Comb-257)

\NewDocumentCommand{\bigrule}{}{%
    \par%
    {%
        \centering%
        \rule{0.6\textwidth}{0.4pt}%
        \par%
    }% typesets a horizontal rule like on the page with the prayers before and after the office.
}  %%If you're using color at all in your document, you might want to either force \myrule to use black or make it cusotmizable. (from u/Independent-Comb-257)


%%% typesets a cross pattée.

\NewDocumentCommand{\mycross}{}{%
{\symbolfont\small{{+}}}%
{}
} %% can replace <+> with <U+2720> if a suitable font is found (or EB Garamond font is fixed); command will be useful in lieu of typing the Unicode.

%% macro to format psalm text

%\NewDocumentCommand{\pstexte}{ m }{%
%    \smallskip%
%    \noindent%
%    \begin{itemize}[%
%    		label=\null, %
%			leftmargin=0pt, %
%			itemindent=0mm, %
%			labelsep=0pt, %
%			labelwidth=0pt, %
%			rightmargin=0pt, %
%			parsep=0pt, %
%			topsep=0pt, %
%			itemsep=0pt]%
%    \input{psaumes/#1}
%    \end{itemize}}
          
      \setlength{\columnsep}{3pc}
%\setlength{\columnsep}{10pt}
\setlength\columnseprule{0.4pt}


%% original version kept so that I can fix errors found later and keep using common texts, going forward
    \NewDocumentCommand{\textebilingue}{ mm }{%
    \smallskip%
       \begin{paracol}{2}
           \input{psaumes/#1}
           \switchcolumn
    \input{psaumes/#2}
    \end{paracol}}

%2nd version;  no need for input (too many chapter/collect files) + is the equivalent of long, allows \par
%    \NewDocumentCommand{\texteparacol}{ +m+m }{%
%    \smallskip%
%       \begin{paracol}{2}
%           #1
%           \switchcolumn
%    #2
%    \end{paracol}}

%%%3rd version: need to switch languages!
    \NewDocumentCommand{\texteparacol}{ +mm+m }{%
    \smallskip%
       \begin{paracol}{2}
           #1
           \switchcolumn
           \begin{otherlanguage}{#2}
    #3
    \end{otherlanguage}
    \end{paracol}}
    
    
    %%psalm vernacular file has language command
    \NewDocumentCommand{\pstexte}{mm}{%
    \smallskip%
    \noindent%%
   \begin{paracol}{2}
    \begin{enumerate}[%
    		label=\arabic*., %
			leftmargin=3pt, %
			itemindent=3mm, %
			labelsep=3pt, %
			labelwidth=0pt, %
			rightmargin=0pt, %
			parsep=0pt, %
			topsep=0pt, %
			itemsep=0pt,%
			start=2]%
    \input{psaumes/#1.tex}
    \end{enumerate}
    \switchcolumn
       \begin{enumerate}[%
    		label=\arabic*.,%
			leftmargin=3pt, %
			itemindent=3mm, %
			labelsep=3pt, %
			labelwidth=0pt, %
			rightmargin=0pt, %
			parsep=0pt, %
			topsep=0pt, %
			itemsep=0pt]%
    \input{psaumes/#2.tex}
      \end{enumerate}
    \end{paracol}%
    } %% modified from original in order to allow 2 column psalms
    
    %% Command de Matthias Bry modifié, prints psalm incipit score with the text
    \NewDocumentCommand{\psalmus}{mmmm}{%
	\needspace{4\baselineskip}%
	\smalltitle{Psalmus #1.}%
	\smallscore[n]{#2}%
	\pstexte{#3}{#4}%
}

    %% sd=semiduplex. especially for Dominica ad Vesperas in the Psalterium (the only example which comes to mind, in fact!)
 \NewDocumentCommand{\sdpsalmusbilingue}{mmm}{%
	\needspace{4\baselineskip}%
		\smalltitle{Psalmus #1.}%
	\pstexte{#2}{#3}%
	}

    %% Command de Matthias Bry modifié, prints canticle incipit score with the text and scriptural commentary
    \NewDocumentCommand{\canticum}{mmmm}{
	\needspace{4\baselineskip}
	\smalltitle{Canticum #1.}
	 \grecommentary{#2}
	\smallscore[n]{#3}
	\pstexte{#4}
}
	
    %% macro to print any additional text (Capitiulum, oratio, rubrics)
 \NewDocumentCommand{\textes}{ m }{%
    \input{textes/#1}}
    
%     \NewDocumentCommand{\lection}{m}{%
%   \begin{multicols}{2}#1\end{multicols}} %% this doesn't work as expected, needs to be fixed
    
    
    
    %% SC work for page headers and for secondary headings but not so much for things like "Dominica…" and certainly not the feast name%%
    
    %%%% Headers%%%%
    \pagestyle{fancy}
\fancyhead{}
\fancyfoot{}
\renewcommand{\headrulewidth}{0pt}


 \setlength{\headheight}{13.59999pt} %% recommended by LaTeX

%\fancyhead[RO]{\small\thepage}
%\fancyhead[LE]{\small\thepage}

\fancyhead[RO]{\small\rightmark\hspace{0.5cm}\thepage}
\fancyhead[LE]{\small\thepage\hspace{0.5cm}\leftmark}

\newcommand{\setheaders}[2]{
	\renewcommand{\rightmark}{{\sc#2}}
	\renewcommand{\leftmark}{{\sc#1}}
}
\setheaders{}{}


%% TITLE Commands %%%% %% TITLE Commands %%%% currently dead but Matthias says this isn't a good idea.
 \titleformat{\chapter}[display]{\Huge\filcenter\sc\addfontfeature{LetterSpace=5.0}}{}{0pt}{} %%
\titleformat{\section}[display]{\LARGE\filcenter\sc\addfontfeature{LetterSpace=5.0}}{}{}{} %%
\titleformat{\subsection}[display]{\Large\filcenter\sc\addfontfeature{LetterSpace=5.0}}{}{}{} %%originally \large
\setcounter{secnumdepth}{0}
\titlespacing*{\chapter}{0pt}{-25pt}{20pt} %%this should reduce spacing before chapter title while putting it flush with the top of the text area
%% see https://tex.stackexchange.com/questions/63390/how-to-decrease-spacing-before-chapter-title
%%\titlespacing*{\chapter}{0pt}{-50pt}{40pt}  %%this puts chapter title in HEADER which we do not want
\titlespacing{\section}{0pt}{*0}{*0}

%%need section titleformat more appropriate for weekdays of psalter and minor feasts 
% See https://tex.stackexchange.com/questions/623797/multiple-section-styles-in-same-document.

\addto\captionsecclesiasticlatin{\renewcommand\contentsname{\centerline{INDEX GENERALIS.}}}

% \NewDocumentCommand{\mysection}{m}{\section{#1}\setheaders{{\scshape\addfontfeature{LetterSpace=5.0}#1}}{{\scshape\addfontfeature{LetterSpace=5.0} #1}}}
 
\NewDocumentCommand{\header}{m}{\setheaders{{\scshape\addfontfeature{LetterSpace=5.0}#1}}{{\scshape\addfontfeature{LetterSpace=5.0} #1}}}

%% all of these are a mess and should probably be ignored, but feel free to use them by referring to Matthias B's originals in the Nocturnale Romanum repository.

\newcommand{\officiumtitulum}[1]{
%  \newpage
%  \thispagestyle{empty}
  \setheaders{{\scshape\addfontfeature{LetterSpace=3.0}#1}}{{\scshape\addfontfeature{LetterSpace=3.0} #1}}
  \begin{center}
  {\scshape\large #1}\par
  \end{center}}
  
%   \setheaders{{{\scshape\addfontfeature{LetterSpace=3.0}}}{{{\scshape\addfontfeature{LetterSpace=3.0} {#1}}}

%%formatting date of feasts where this listed at top of page
\NewDocumentCommand{\litdate}{m}{%
\smalltitle{Die #1.}%
}

\NewDocumentCommand{\litdateen}{m}{%
\smalltitle{#1.}%
}

%%formatting for Sunday feasts : inter, post… name is for consistency
\NewDocumentCommand{\dominicadate}{m}{%
\smalltitle{Dominica #1.}%
}

\NewDocumentCommand{\sundaydate}{m}{%
\smalltitle{Sunday #1.}%
}

%%will use litdaten for English
%%formatting for feasts that fall on a weekday fixed to some other thing, like Corpus Christi, Sacred Heart, BVM in Passiontide, and the new feasts found in the appendix (Feria Roman numeral in Latin: inter, post… name is for consistency
\NewDocumentCommand{\feriadate}{m}{%
\smalltitle{Feria #1.}%
}

 %%looks like xparse formatting changed so o just has No Default Value whereas O{} is that OR you can insert something in braces!
%%answer thanks to David Carlisle
\NewDocumentCommand{\duplexclassis}{mo}{%
\needspace{5\baselineskip}%
\vspace{\baselineskip}%
 {\centering\textit{\large Duplex #1 classis\IfValueT{#2}{ #2}.}\par}%
}

\NewDocumentCommand{\doubleclass}{mo}{%
\needspace{5\baselineskip}%
\vspace{\baselineskip}%
 {\centering\textit{\large Double of the #1 class\IfValueT{#2}{ #2}.}\par}%
}


\NewDocumentCommand{\duplexmajus}{}{%
\needspace{5\baselineskip}%
\vspace{\baselineskip}%
 {\centering{\textit{\large Duplex majus.}}\par}%
}

\NewDocumentCommand{\doublemajor}{}{%
\needspace{5\baselineskip}%
\vspace{\baselineskip}%
 {\centering{\textit{\large Double major.}}\par}%
}

%%is Duplex even specified.as it's the default?
\NewDocumentCommand{\duplex}{}{%
\needspace{5\baselineskip}%
\vspace{\baselineskip}%
 {\centering{\textit{\large Duplex.}}\par}%
}

\NewDocumentCommand{\double}{}{%
\needspace{5\baselineskip}%
\vspace{\baselineskip}%
 {\centering{\textit{\large Double.}}\par}%
}

\NewDocumentCommand{\semiduplex}{}{%
\needspace{5\baselineskip}%
\vspace{\baselineskip}%
 {\centering{\textit{\large semiuplex.}}\par}%
}

\NewDocumentCommand{\semidouble}{}{%
\needspace{5\baselineskip}%
\vspace{\baselineskip}%
 {\centering{\textit{\large Semidouble.}}\par}%
}

\NewDocumentCommand{\simplex}{}{%
\needspace{5\baselineskip}%
\vspace{\baselineskip}%
 {\centering{\textit{\large Simplex.}}\par}%
}

\NewDocumentCommand{\simple}{}{%
\needspace{5\baselineskip}%
\vspace{\baselineskip}%
 {\centering{\textit{\large Simple.}}\par}%
}

\NewDocumentCommand{\rank}{m}{
\needspace{5\baselineskip}  %% very small if there are neumes above the staff, including flats in mode 2, e.g. O Doctor optime
\vspace{\baselineskip}
 {\centering{\textit{\large #1}}\par}
}



\NewDocumentCommand{\bigtitle}{m}{
\needspace{5\baselineskip}  %% very small if there are neumes above the staff, including flats in mode 2, e.g. O Doctor optime
\vspace{\baselineskip}
 {\centering{\large#1}\par}
}

\NewDocumentCommand{\smalltitle}{m}{
\needspace{5\baselineskip}  %% very small if there are neumes above the staff, including flats in mode 2, e.g. O Doctor optime
\vspace{\baselineskip}
 {\centering #1\par}
}

\NewDocumentCommand{\chapterlateng}{}{%
\smalltitle{Chapter.}%
}


%\newcommand{\subtitulum}[1]{
% \subsection{
%  \begin{center}
%  {\large#1}\par
%  \end{center}}}
  
  %% something has broken here
  %% see https://tex.stackexchange.com/questions/545841/paragraph-ended-before-ttlformatsi-was-complete

%% originally used apparently for typesetting portions of the Common of the BVM
\newcommand{\espacetitre}{\vspace{-0.5cm}}
\newcommand{\espaceps}{\vspace{-3mm}} %% ps =psaume
\newcommand{\espacecap}{\vspace{-3mm}} %% cap is capitulum
\newcommand{\espace}{\vspace{2mm}}

\NewDocumentCommand{\secreto}{}{Pater noster. \rubrique{secreto.}}

\NewDocumentCommand{\pateravecredo}{}{\large{Pater noster, Ave María, \textit{et} Credo.}}

\NewDocumentCommand{\pateravedeus}{}{%
Pater noster et Ave María. ℣. Deus in adjutórium.%
}

\NewDocumentCommand{\prayers}{}{%
Our Father and Hail Mary, \rubrique{silently.}%
}

\NewDocumentCommand{\orationes}{}{%
Pater noster et Ave María \rubrique{secreto.}%
}

%%for litanies
\makeatletter  
\newcounter{score}
\newcounter{tabstop}[score]
\newcommand{\grealign}{%
	\@bsphack%
	\ifgre@boxing\else%
		\kern\gre@dimen@begindifference%
		\stepcounter{tabstop}%
		\expandafter\zsavepos{stop-\thescore-\thetabstop}%
		\kern-\gre@dimen@begindifference%
	\fi%
	\@esphack%
}

\newcommand{\setstops}{%
  \gdef\nstabbing@stops{%
    \hspace*{-\oddsidemargin}\hspace{-1in}%
    \hspace*{\zposx{stop-\thescore-1} sp}\=%
  }%
  \count@=\@ne
  \loop\ifnum\count@<\value{tabstop}%
    \begingroup\edef\x{\endgroup
      \noexpand\g@addto@macro\noexpand\nstabbing@stops{%
        \noexpand\hspace{-\noexpand\zposx{stop-\thescore-\the\count@} sp}%
        \noexpand\hspace{\noexpand\zposx{stop-\thescore-\the\numexpr\count@+1} sp}\noexpand\=%
      }%
    }\x
    \advance\count@\@ne
  \repeat
  \nstabbing@stops\kill
}
\makeatother

\newenvironment{nstabbing}
  {\setlength{\topsep}{0pt}%
   \setlength{\partopsep}{0pt}%
   \tabbing%
   \setstops}
  {\endtabbing\stepcounter{score}}

\makeatletter
\newcounter{blankpages} %% should remove number on blank page at end of preface. Cf https://tex.stackexchange.com/questions/250761/how-to-not-count-blank-pages-in-frontmatter-of-two-sided-document
\def\cleardoublepage{%
    \clearpage
    \if@twoside
    \ifodd\c@page
    \else
    \hbox{}
    \thispagestyle{empty}%
    \newpage\stepcounter{blankpages}%
    \if@twocolumn\hbox{}\newpage\fi
    \fi
    \fi
}
\newcommand{\@romannoblank}[1]{%
    \@roman{\numexpr#1-\value{blankpages}\relax}%
}
\makeatother

\providecommand\phantomsection{} %% should make hyperref work properly %%only works with section etc., not with label etc. (that requires \phantomsection
%

\NewDocumentEnvironment{enpars}{}
  {\begin{otherlanguage*}{english}}
  {\par\end{otherlanguage*}}
  
  \NewDocumentEnvironment{frpars}{}
  {\begin{otherlanguage*}{french}}
  {\par\end{otherlanguage*}}
%\renewcommand{\contentsname}{\hfill\bfseries\LARGE Index Generalis.\hfill\null}  %% modifies TOC title
%\renewcommand{\cfttoctitlefont}{\hfil\LARGE}

%%%SUBFILES%%%
\usepackage{xr} %%to allow cross-references between documents%%%
  \usepackage{subfiles}
  
  %% When we start a new subfile (new chapter or section) , 
%% we start on a new page (with blank filling on the previous page) and create a corresponding label.

\newcommand{\customsubfile}[1]{\newpage\label{#1}\thispagestyle{empty}\subfile{#1}}
 %%% will need to fix spacing problem in psalms at a later date.
\begin{document}
\thispagestyle{empty}
\smalltitle{The Major Litanies on the feast of Saint Mark and the Minor Litanies before the Ascension}
%\addcontentsline{toc}{section}{For processions.} \phantomsection

{\centering{and in other processions.\footnote{The special invocations of the litany and the prayers for other occasions are found in the ritual.}\par}}

\rubrique{Ante Processionem, cantatur stando. Before the procession, the following is sung, standing.}

\gscore[Ant.]{2.}{an_exsurge_domine_adjuva_en_solesmes}

\rubrique{Duo cantores ante altare majus genuflexi Litanias incipiunt, et omnes versus duplicantur, nisi Processio fuerit impedita. Surgunt omnes et ordinatim procendunt, ad
\normaltext{Sancta Dei}.}

\rubrique{Two cantors kneeling before the main altar begin the Litany. All verses are doubled, except if the procession is impeded. All rise and the procession departs at \normaltext{Sancta Dei.}}

\rubrique{\normaltext{Litaniæ Sanctorum, \pageref{litaniae_sanctorum},} ps.\@ 69 et preces.}

\rubrique{\normaltext{Litany of the Saints, \pageref{litaniae_sanctorum}.} with ps.\@ 69 and the prayers.}

\newpage

\smalltitle{Forty Hours'\thinspace Devotion.}

\addcontentsline{toc}{section}{Forty Hours'\thinspace Devotion}\phantomsection
%. After the Mass on the day of exposition.
\texteparacol{In die expositionis, missa finita, Celebrans ad scammum deponit planetam et
manipulum et assumit pluviale. S. Ministri deponunt manipula.
 Celebrans ad scammum incensum in duo thuribula sine benedictione injicit.
 
Celebrans genuflexus in infimo gradu Altaris SS.\@ Sacramentum incensat. Deinde assumit velum humerale.
 
Fit processio. Intonatur \textit{Pange lingua,} ut infra.
 
Peracta processione, Diaconus SS. Sacramentum in throno collocat, interim Celebrans deponit velum humerale, et deinde intonatur
\textit{Tantum ergo Sacramentum.} Ad \textit{Genitori} incensatur SS.\@ Sacramentum
more solito.

Completo cantu hymni \textit{Tantum ergo,} non additur \textit{Panem de cælo,}
etc., sed statim Cantores incipiunt Litanias Sanctorum, \normaltext{\pageref{litaniae_sanctorum}.}}
{english}
{After Mass on the day of exposition, the celebrant removes the chasuble and maniple
at the sedilia and puts on the cope. The sacred ministers remove their maniples. The celebrant imposes incense without blessing it in the two thuribles.

The celebrant, kneeling on the lowest grade of the altar steps, incenses the Most Blessed Sacrament, then he puts on the humeral veil.

The procession takes place. \textit{Pangue lingua} is intoned, as below.

At the end of the procession, the deacon places the Most Blessed Sacrament on the throne, while the celebrant removes the humeral veil. Then \textit{Tantum ergo Sacramentum} is intoned. At \textit{Genitori,} the Most Blessed Sacrament is incensed in the usual way.

After the hymn \textit{Panem de cælo} is not said, but the cantors immediately begin the Litany of the Saints, \normaltext{\pageref{litaniae_sanctorum}.}}

\label{hy_pange_lingua_corpus_christi_solesmes} \phantomsection
\gscore[]{3.}{hy_pange_lingua_corpus_christi_solesmes}
\textes{pange_lingua_gloriosi_en}

\myrule
\rubrique{In die repositionis tantum. On the day of reposition only.}
\texteparacol{\vv Panem de cælo præstitísti eis. \tpalleluia{}\footnote{Atque per octavam Corporis Christi. \textit{As well as throughout the octave of Corpus Christi.}}

\rr Omne dele\-ctaméntum in se habéntem. \tpalleluia{}}
{english}
{\noindent \vv Thou hast given them bread from heaven. (\textit{P.T.} Alleluia.)

\noindent \rr Having within it all sweetness. (\textit{P.T.} Alleluia.)}

\myrule

\textes{xl_horarum_orationes}

\texteparacol{In die repositionis, post Missam Celebrans, ad scammum, deponit casulam et
manipulum et assumit pluviale. S. Ministri manipula deponunt, et omnes ante Altare redeunt.

Cantantur Litaniæ Sanctorum et preces usque ad ℣. \textit{Dómine exáudi oratiónem meam}
inclusive.

Surgit Celebrans et incensum in duo thuribula sine benedictione injicit, SS.\@ Sacramentum incensat, et assumit velum humerale.

Fit processio. Intonatur \textit{Pange lingua,} \pageref{hy_pange_lingua_corpus_christi_solesmes}.

Peracta processione, et SS.\@ Sacramento super Altare collocato, Celebrans deponit velum humerale.

Intonatur
\textit{Tantum ergo.} Ad \textit{Genitori,} Celebrans incemsum injicit et incensatur SS.\@ Sacramentum.

Completo cantu hymni \textit{Tantum ergo,} intonatur \textit{Panem de cælo,} etc.

Surgit Celebrans et, quin præmittat \textit{Dóminus vobíscum,} cantat omnes orationes præscri\-ptas.

Quibus absolutis rursus genuflectit et cantat \textit{Dómine exáudi} etc. Deinde cantores cantant \textit{Exáudiat nos,} etc., et Celebrans submissa voce addit \textit{Et fidélium,} etc.

Celebrans assumit velum humerale, et Benedictionem impertit more solito.}
{english}
{After Mass on the day of reposition, the celebrant removes the chasuble and maniple
at the sedilia and puts on the cope. The sacred ministers remove their maniples.

The Litany of the Saints, is sung with its prayers until ℣.
\textit{O Lord, hear my prayer} inclusively, \normaltext{\pageref{litaniae_sanctorum},}

The celebrant rises and imposes incense without blessing it in the two thuribles, and then he puts on the humeral veil.

The procession takes place. \textit{Pangue lingua} is intoned, \normaltext{\pageref{hy_pange_lingua_corpus_christi_solesmes}.}

At the end of the procession, the Most Blessed Sacrament having been replaced on the throne, the celebrant removes the humeral veil. Then \textit{Tantum ergo Sacramentum} is intoned. At \textit{Genitori,} the celebrant imposes incense and Most Blessed Sacrament is incensed in the usual way.

Having finished the hymn \textit{Tantum ergo,} \textit{Panem de cælo,} etc. are sung.

The celebrant rises, and without saying \textit{Dóminus vobíscum,} sings the prescribed prayers.

Having finished these, he kneels again and sings \textit{Dómine exáudi} etc. Then the cantors sing
\textit{Exáudiat nos,} etc., and the celebrant adds in a lower voice. \textit{Et fidélium,} etc.

The celebrant puts on the humeral veil again, and Benediction is given in the usual way.}

%\newpage
%
%\newpage

%%litanies align with odd page as first one

\cleardoublepage

\label{litaniae_sanctorum}
\smalltitle{Litany of the Saints.\footnote{Translation of the orations © Michael Martin. Used with permission.}}

\smallscore{litaniae_sanctorum_1_kyrie_eleison_solesmes}

\gresetlastline{justified}
\smallscore{litaniae_sanctorum_2_pater_solesmes}%
\begin{nstabbing}
\>Fili \>Re-\>dém-\>ptor mundi \>\>\textbf{De}-\>us, \>mi-\>se-\>\textit{ré}-\>\textit{re} \>no-\>bis.\\
%\>Fili \>Re-\>démp-\>tor \>mundi \>\textbf{De}-\>us, \>mi-\>se-\>\textit{ré}-\>\textit{re} \>no-\>bis.\\
\>Spí-\>ri-\>tus \>Sancte \>\>\textbf{De}-\>us, \>mi-\>se-\>\textit{ré}-\>\textit{re} \>no-\>bis.\\
\>Sancta \>\>Trínitas \>\>unus \>\textbf{De}-\>us, \>mi-\>se-\>\textit{ré}-\>\textit{re} \>no-\>bis.\\
\end{nstabbing}

%%can we improve alignment of dactyls?
%%2nd syllables are all misaligned
%%only vir- is aligned
%%last are also wrong, but it's subtly in between e and l for Angels
%%however it resembles example http://gregorio-project.github.io/tips/litany.html
%%so for now we will call it a day
\smallscore{litaniae_sanctorum_3_sancta_maria_solesmes}
\begin{nstabbing}
\>Sancta Dé-\>\>i \>\textbf{Gé}-\>ni- \>trix, \>ó-\>\textit{ra} \>\textit{pro} \>nó-\>bis.\\
\>Sancta Vír-\>\>go \>\textbf{vír-}\>gi-\>num, \>ó-\>\textit{ra} \>\textit{pro} \>nó-\>bis.\\
\>Sancte \>\>\>\textbf{Mí-}\>cha-\>el, \>ó-\>\textit{ra} \>\textit{pro} \>nó-\>bis.\\
\>Sancte \>\>\>\textbf{Gá-}\>bri-\>el, \>ó-\>\textit{ra} \>\textit{pro} \>nó-\>bis.\\
\>Sancte \>\>\>\textbf{Rá-}\>pha-\>el, \>ó-\>\textit{ra} \>\textit{pro} \>nó-\>bis.\\
\end{nstabbing}

\smallscore{litaniae_sanctorum_4_omnes_solesmes}
%\noindent
\begin{nstabbing}
\>Omnes \quad sancti \quad beatórum Spirítuum \>\>\>\>\>\textbf{ór}di-nes, \>\>o-\>rá-\>\textit{te} \>\textit{pro} \>nó-\>bis.\\
\>Sancte \>Joánnes \>\>Ba\>\textbf{ptí}-\>\>sta, \>\>o-\>\textit{ra} \>\textit{pro} \>nó-\>bis.\\
\>Sancte \>\>\>\>\>\textbf{Jó-}\>seph, \>\>o-\>\textit{ra} \>\textit{pro} \>nó-\>bis.\\
\>Omnes \>sancti \>Patriárchæ et \quad \>Pro-\>\>\textbf{phé}-\>tæ, \>o-\>rá-\>\textit{te} \>\textit{pro} \>nó-\>bis.\\
\end{nstabbing}
\gresetlastline{ragged}
\smallscore{litaniae_sanctorum_5_sancte_petre_solesmes}
\begin{multicols}{2}
\noindent
San\-cte \textbf{Pau}le,\hfill o\textit{ra}.\\
San\-cte An\textbf{dré}a,\hfill o\textit{ra}.\\
San\-cte Ja\textbf{có}be,\hfill o\textit{ra}.\\
San\-cte Jo\textbf{án}nes,\hfill o\textit{ra}.\\
San\-cte \textbf{Tho}ma,\hfill o\textit{ra}.\\
San\-cte Ja\textbf{có}be,\hfill o\textit{ra}.\\
San\-cte Phi\textbf{líp}pe,\hfill o\textit{ra}.\\
San\-cte Bartholo\textbf{mǽ}e,\hfill o\textit{ra}.\\
San\-cte Mat\textbf{thǽ}e,\hfill o\textit{ra}.\\
San\-cte \textbf{Si}mon,\hfill o\textit{ra}.\\
San\-cte Thad\textbf{dǽ}e,\hfill o\textit{ra}.\\
San\-cte Mat\textbf{thí}a,\hfill o\textit{ra}.\\
San\-cte \textbf{Bár}naba,\hfill o\textit{ra}.\\
San\-cte \textbf{Lu}ca,\hfill o\textit{ra}.\\
San\-cte \textbf{Mar}ce,\hfill o\textit{ra}.\\
Omnes san\-cti Apóstoli et \\ \emspace Evange\textbf{lí}stæ,\hfill orá\textit{te}.\\
Omnes san\-cti Discípuli \textbf{Dó}mini,\hfill orá\textit{te}.\\
Omnes san\-cti Inno\textbf{cén}tes,\hfill orá\textit{te}.\\
San\-cte \textbf{Sté}phane,\hfill o\textit{ra}.\\
San\-cte Lau\textbf{rén}ti,\hfill o\textit{ra}.\\
San\-cte Vin\textbf{cén}ti,\hfill o\textit{ra}.\\
San\-cti Fabiáne et Sebasti\textbf{án}e,\hfill orá\textit{te}.\\
San\-cti Joánnes et \textbf{Pau}le,\hfill orá\textit{te}.\\
San\-cti Cosma et Dami\textbf{án}e,\hfill orá\textit{te}.\\
San\-cti Gervási et Pro\textbf{tá}si,\hfill orá\textit{te}.\\
Omnes san\-cti \textbf{Már}tyres,\hfill orá\textit{te}.\\
San\-cte Sil\textbf{vé}ster,\hfill o\textit{ra}.\\
San\-cte Gre\textbf{gó}ri,\hfill o\textit{ra}.\\
San\-cte Am\textbf{bró}si,\hfill o\textit{ra}.\\
San\-cte Augus\textbf{tí}ne,\hfill o\textit{ra}.\\
San\-cte Hie\textbf{ró}nyme,\hfill o\textit{ra}.\\
San\-cte Mar\textbf{tí}ne,\hfill o\textit{ra}.\\
San\-cte Nico\textbf{lá}e,\hfill o\textit{ra}.\\
Omnes san\-cti Pontífices et \\ \emspace Confes\textbf{só}res,\hfill orá\textit{te}.\\
Omnes san\-cti Do\textbf{ctó}res,\hfill orá\textit{te}.\\
San\-cte An\textbf{tó}ni,\hfill o\textit{ra}.\\
San\-cte Bene\textbf{dí}cte,\hfill o\textit{ra}.\\
San\-cte Ber\textbf{nár}de,\hfill o\textit{ra}.\\
San\-cte Do\textbf{mí}nice,\hfill o\textit{ra}.\\
San\-cte Fran\textbf{cí}sce,\hfill o\textit{ra}.\\
Omnes san\-cti Sacerdótes et \\  \emspace Le\textbf{ví}tæ,\hfill orá\textit{te}.\\
Omnes san\-cti Mónachi et \\ \emspace Ere\textbf{mí}tæ,\hfill orá\textit{te}.\\
San\-cta María Magda\textbf{lé}na,\hfill o\textit{ra}.\\
San\-cta \textbf{A}gatha,\hfill o\textit{ra}.\\
San\-cta \textbf{Lú}cia,\hfill o\textit{ra}.\\
San\-cta \textbf{A}gnes,\hfill o\textit{ra}.\\
San\-cta Cæ\textbf{cí}lia,\hfill o\textit{ra}.\\
San\-cta Catha\textbf{rí}na,\hfill o\textit{ra}.\\
San\-cta Anas\textbf{tá}sia,\hfill o\textit{ra}.\\
Omnes san\-ctæ Vírgines \\ \emspace et \textbf{Ví}duæ,
\hfill orá\textit{te}.\\
Omnes San\-cti et San\-ctæ \textbf{De}i,\\ 
\zeroptspace \hfill intercédi\textit{te pro} nobis.
\end {multicols}

%%less than ideal formatting shared with gregobase
\smallscore{litaniae_sanctorum_6_propitius_esto_solesmes}%

\noindent
Propí\textit{tius} \textbf{es}to,\hfill exáudi nos Dómine.\\
Ab \textit{omni} \textbf{ma}lo,\hfill líbera nos Dómine.\\
Ab \textit{omni pec}\textbf{cá}to,\hfill líbera nos Dómine.\\
Ab \textit{ira} \textbf{tu}a,\hfill líbera nos Dómine.\\
A subitánea et impro\textit{vísa} \textbf{mor}te,\hfill líbera nos Dómine.\\
Ab insídi\textit{is di}\textbf{á}boli,\hfill líbera nos Dómine.\\
Ab ira, et ódio, et omni mala \textit{volun}\textbf{tá}te,\hfill líbera nos Dómine.\\
A spíritu forni\textit{cati}\textbf{ó}nis,\hfill líbera nos Dómine.\\
A fúlgure et \textit{tempes}\textbf{tá}te,\hfill líbera nos Dómine.\\
A flagéllo \textit{terræ}\textbf{mó}tus,\hfill líbera nos Dómine.\\
A peste, fa\textit{me, et} \textbf{bel}lo,\hfill líbera nos Dómine.

\myrule

\noindent Ab imminénti\textit{bus} \textit{pér}iculis,\footnote{\rubrique{In Oratione Quadraginta Horarum tantum. At the Forty Hours' \thinspace Prayer only.}} \hfill líbera nos Dómine.
%\negativebaselineskipvspace
\myrule
\noindent A mor\textit{te per}\textbf{pé}tua,\hfill líbera nos Dómine.\\
Per mystérium sanctæ incarnati\tinyhspace\textit{ónis} \textbf{tu}æ,\hfill líbera nos Dómine.\\
Per ad\textit{véntum} \textbf{tu}um,\hfill líbera nos Dómine.\\
Per nativi\tinyhspace\textit{tátem} \textbf{tu}am,\hfill líbera nos Dómine.\\
Per baptísmum et sanctum jejú\textit{nium} \textbf{tu}um,\hfill líbera nos Dómine.\\
Per crucem et passi\tinyhspace\textit{ónem} \textbf{tu}am,\hfill líbera nos Dómine.\\
Per mortem et sepul\tinyhspace\textit{túram} \textbf{tu}am,\hfill líbera nos Dómine.\\
Per sanctam resurrecti\tinyhspace\textit{ónem} \textbf{tu}am,\hfill líbera nos Dómine.\\
Per admirábilem ascensi\tinyhspace\textit{ónem} \textbf{tu}am,\hfill líbera nos Dómine.\\
Per advéntum Spíritus San\tinyhspace\textit{cti Pa}\textbf{rá}cliti,\hfill líbera nos Dómine.\\
In di\tinyhspace\textit{e ju}\textbf{dí}cii,\hfill líbera nos Dómine.
%\end{nstabbing}

\smallscore{litaniae_sanctorum_7_peccatores_solesmes}%
%\begin{nstabbing}
%Ut nó\textit{bis} \textbf{pár}cas, te rógamus áudi nos.\\
%Ut Ecclésiam túam san\-ctam ' régere et conserváre \textit{di}\textbf{gne}ris,\\
%\hspace*{1cm} te rógamus áudi nos.
%\end{nstabbing}
\noindent
Ut no\textit{bis} \textbf{par}cas,\hfill te rogámus audi nos.\\
Ut nobis \textit{in}\textbf{dúl}geas,\hfill te rogámus audi nos.\\
Ut ad veram pæniténtiam nos perdúcere
\textit{di}\textbf{gné}ris,\hfill te rogámus audi nos.\\
Ut Ecclésiam tuam sanctam ’ régere et
conserváre\\
\emspace \textit{di}\textbf{gné}ris,\hfill te rogámus audi nos.\\
Ut Domnum Apostólicum et omnes
ecclesiásticos \\
\emspace órdines ’ in sancta  religióne conserváre \textit{di}\textbf{gné}ris,\hfill te rogámus audi nos.\\
Ut inimícos sanctæ Ecclésiæ ’ humiliáre
\textit{di}\textbf{gné}ris,\hfill te rogámus audi nos.\\
Ut régibus et princípibus christiánis ’ pacem et \\
\emspace veram concórdiam donáre \textit{di}\textbf{gné}ris,\hfill te rogámus audi nos.\\
Ut cuncto pópulo christiáno ’ pacem et
unitátem largíri \textit{di}\textbf{gné}ris,\hfill te rogámus audi nos.\\
Ut omnes errántes ad unitátem Ecclésiæ \\
\emspace
revocáre, ’ et infidéles univérsos ad Evangélii
lumen \\
\emspace perdúcere \textit{di}\textbf{gné}ris,\hfill te rogámus audi nos.\\
Ut nosmetípsos in tuo sancto servítio ’
confortáre \\
\emspace  et conserváre \textit{di}\textbf{gné}ris,\hfill te rogámus audi nos.\\
Ut mentes nostras ’ ad cæléstia desidéri\tinyhspace\textit{a}
\textbf{é}rigas,\hfill te rogámus audi nos.\\
Ut ómnibus benefactóribus nostris ’
sempitérna bona \textit{re}\textbf{trí}buas,\hfill te rogámus audi nos.\\
Ut ánimas nostras ’ fratrum, propinquórum et \\
\emspace benefactórum nostrórum ’ ab ætérna
damnatióne \textit{e}\textbf{rí}pias,\hfill te rogámus audi nos.\\
Ut fructus terræ ’ dare et conserváre
\textit{di}\textbf{gné}ris,\hfill te rogámus audi nos.\\
Ut ómnibus fidélibus defúnctis ’ réquiem
ætérnam \\
\emspace donáre \textit{di}\textbf{gné}ris,\hfill te rogámus audi nos.\\
Ut nos exaudíre \textit{di}\textbf{gné}ris,\hfill te rogámus audi nos.\\
Fi\tinyhspace\textit{li} \textbf{De}i,\hfill te rogámus audi nos.

\smallscore{litaniae_sanctorum_8_agnus_Dei_solesmes}

\smallscore{litaniae_sanctorum_9_kyrie_eleison_english_solesmes}

\textes{litaniae_san_ps_vv}

\textes{litaniae_sanctorum_orationes}
\end{document}